% Standalone problem document, generated from the adjacent README.md.
% Compile with XeLaTeX. Mathematical content is maintained in Markdown.
\documentclass[11pt,a4paper]{article}
\usepackage[fontsize=10.5pt]{scrextend}
\usepackage[margin=22mm,headheight=15pt,headsep=9mm,footskip=11mm]{geometry}
\usepackage{fontspec}
\setmainfont{texgyrepagella}[Extension=.otf,UprightFont=*-regular,BoldFont=*-bold,ItalicFont=*-italic,BoldItalicFont=*-bolditalic]
\setsansfont{texgyreheros}[Extension=.otf,UprightFont=*-regular,BoldFont=*-bold,ItalicFont=*-italic,BoldItalicFont=*-bolditalic]
\setmonofont{texgyrecursor}[Extension=.otf,UprightFont=*-regular,BoldFont=*-bold,ItalicFont=*-italic,BoldItalicFont=*-bolditalic,Scale=MatchLowercase]
\usepackage{amsmath,amssymb}
\usepackage{unicode-math}
\setmathfont{texgyrepagella-math.otf}
\usepackage{xcolor,microtype,enumitem,needspace,fancyhdr}
\usepackage{bookmark}
\definecolor{ink}{HTML}{17344B}
\definecolor{accent}{HTML}{197D80}
\definecolor{muted}{HTML}{596773}
\hypersetup{colorlinks=true,linkcolor=accent,urlcolor=accent,pdftitle={RA-20 - Critical-point
counts for symmetric rank-two approximation with diagonal
zeros},pdfauthor={Open Problems in Numerical Linear Algebra}}
\urlstyle{same}
\setlength{\parindent}{0pt}
\setlength{\parskip}{5pt plus 1pt minus 1pt}
\linespread{1.04}
\setlength{\emergencystretch}{3em}
\widowpenalty=10000
\clubpenalty=10000
\displaywidowpenalty=10000
\allowdisplaybreaks[1]
\setlist{nosep,leftmargin=1.5em}
\providecommand{\tightlist}{\setlength{\itemsep}{0pt}\setlength{\parskip}{0pt}}
\providecommand{\pandocbounded}[1]{#1}
\pagestyle{fancy}
\fancyhf{}
\fancyhead[L]{\sffamily\footnotesize\color{muted}OPEN PROBLEMS IN NUMERICAL LINEAR ALGEBRA}
\fancyhead[R]{\sffamily\bfseries\footnotesize\color{accent}RA-20}
\fancyfoot[L]{\parbox[b]{0.9\textwidth}{\sffamily\fontsize{8}{10}\selectfont\color{muted}Editorial ratings; a bounded search cannot certify that a problem remains unsolved.\\Literature check: 2026-09-11}}
\fancyfoot[R]{\sffamily\footnotesize\color{muted}\thepage}
\renewcommand{\headrulewidth}{0.3pt}
\renewcommand{\headrule}{\hbox to\headwidth{\color{accent}\leaders\hrule height \headrulewidth\hfill}}
\makeatletter
\renewcommand{\section}{\Needspace{5\baselineskip}\@startsection{section}{1}{0pt}{12pt plus 2pt minus 2pt}{4pt}{\sffamily\bfseries\large\color{ink}}}
\renewcommand{\subsection}{\Needspace{4\baselineskip}\@startsection{subsection}{2}{0pt}{9pt plus 1pt minus 1pt}{3pt}{\sffamily\bfseries\normalsize\color{ink}}}
\makeatother
\setcounter{secnumdepth}{0}
\begin{document}
{\sffamily\small\color{muted}Randomized and low-rank approximation\par}
\vspace{3pt}
{\raggedright\hyphenpenalty=10000\exhyphenpenalty=10000\sffamily\bfseries\fontsize{21}{24}\selectfont\color{ink}Critical-point
counts for symmetric rank-two approximation with diagonal zeros\par}
\vspace{7pt}
{\sffamily\small\textbf{Difficulty:} challenging\quad\textbf{Importance:} interesting
to specialist\par}
{\sffamily\small\bfseries\color{accent}Status: Open\par}
\vspace{4pt}
{\color{accent}\hrule height 0.8pt}
\vspace{6pt}
\textbf{Topic:} Symmetric structured low-rank approximation

\textbf{Rating rationale:} Four linked formulas predict how diagonal
constraints alter the algebraic complexity of rank-two approximation. A
proof must handle the geometry of the constrained symmetric varieties
uniformly in matrix order.

\subsection{Statement}\label{statement}

For \(s\in\{1,2,3,4\}\) and \(n\ge\max\{3,s\}\), define

\[W_{n,s}=\{X\in\mathbb C^{n\times n}:X^{\mathsf T}=X,
\ \operatorname{rank}X\le2,\ x_{11}=\cdots=x_{ss}=0\}.\]

For generic symmetric \(U\in\mathbb C^{n\times n}\), let \(e_{n,s}\) be
the number of complex critical points on the smooth locus of \(W_{n,s}\)
of

\[d_U(X)=\sum_i(x_{ii}-u_{ii})^2+
2\sum_{i<j}(x_{ij}-u_{ij})^2.\]

Generic means outside a proper algebraic exceptional set. A point is
critical if the differential vanishes on its tangent space. Thus
\(e_{n,s}\) is the Euclidean distance degree for the bilinear extension
of the \textbf{full Frobenius metric}; the off-diagonal terms have
weight two, and complex conjugation is absent.

\textbf{Conjecture (Kubjas--Sodomaco--Tsigaridas, Conjecture 5.6,
\(n\ge3\)).}

\[e_{n,s}=\begin{cases}
3(n-1)-2,&s=1,\\
9(n-2)-2,&s=2,\\
27(n-3)+4,&s=3,\\
81(n-4)+28,&s=4.
\end{cases}\]

All four zero counts form one target. The explicit \(n\ge3\) restriction
avoids a degenerate endpoint in the printed statement: when \(n=2\), the
rank bound is vacuous and both possible varieties are linear with ED
degree one, whereas the displayed \(s=2\) formula does not apply.

\subsection{Evidence and numerical
significance}\label{evidence-and-numerical-significance}

The source's Table 7 gives matching computations through order ten.
These do not constitute an all-dimension proof. The problem counts
stationary candidates for symmetric Frobenius approximation with
prescribed diagonal zeros. It concerns fixed rank two, unlike
\href{https://github.com/ajt60gaibb/OpenProblemsInNLA/blob/main/randomized-and-low-rank-approximation/RA-19/README.md}{corank-one
approximation in general square matrices}.

\newpage
\subsection{References and status
check}\label{references-and-status-check}

\begin{itemize}
\tightlist
\item
  K. Kubjas, L. Sodomaco and E. Tsigaridas, \emph{Exact solutions in
  low-rank approximation with zeros}, Linear Algebra and its
  Applications 641 (2022), 67--97.
  \href{https://doi.org/10.1016/j.laa.2022.01.021}{DOI};
  \href{https://arxiv.org/abs/2010.15636v2}{current author manuscript},
  29 January 2022. Conjecture 5.6 and Table 7, manuscript p.21; §2 and
  §5 supply the distance convention.
\end{itemize}

On 2026-09-11, checked arXiv v2, publication lists and targeted searches
for the paper title, symmetric zero patterns, Conjecture 5.6, Euclidean
distance degree, and later proofs or counterexamples. No full resolution
in the displayed range was found. Source formulas were compared with
Table 7; the restriction excluding \(n=2\) is stated above. The separate
nonsymmetric formulas in Conjecture 5.2 have a table/label discrepancy
and are not imported here. This is a bounded status check.
\end{document}
